\documentclass{article}
\usepackage[T1]{fontenc}
\usepackage{lmodern}
\usepackage{graphicx}
\usepackage{amsmath,amssymb}
\usepackage[a4paper,margin=2cm]{geometry}
\usepackage[absolute,overlay]{textpos}
\setlength{\TPHorizModule}{1mm}
\setlength{\TPVertModule}{1mm}
\usepackage[indonesian]{babel}
\usepackage{hyperref}
\setlength{\emergencystretch}{2em}
% unit: Halaman 70

\begin{document}

% === Halaman 70 (python/native) ===

70

\% Nama file: LSB\_ekstrak.m

\% Ekstraksi pesan dari citra stegano dengan metode modifikasi LSB. Panjang

\% pesan (L) telah disimpan di dalam citra stegano. 

\% Asumsi: jumlah bit untuk merepresentasikan panjang pesan maksimal 30 bit.

\% Input: Citra stegano (berwarna)

\% Output: pesan (teks)

\% Diprogram oleh Rinaldi Munir, @Informatika STEI-ITB


clear all;


\% Baca citra stegano

namaFile = input('Ketikkan nama file citra stego: ', 's');

stegoImage = imread(namaFile);

imshow(namaFile); title ('Citra stego ');


\% Ekstraksi bit-bit panjang pesan terlebih dahulu 

\% Diperlukan 30 byte = 10 pixel x 3 byte pada baris pertama citra stegano

idx = 1;

for col = 1:10

    for i = 1:3

        nL(idx) = bitget(stegoImage(1,col, i), 1);

        idx = idx + 1;

    end    

end

B. Program ekstraksi pesan

\end{document}
